\section{An exact dual value without an attaining multiplier}\label{app:attainment}
Example~13 of the December 15, 2025 manuscript of
\citet{LefebvreSchmidt2025} is stated to admit no finite exact penalty
representation. Under its supremum-based Definition~1, however, the
example has an exact dual value at every nonnegative penalty.
In renamed variables its native set is
\[
 X=\{(u,s,z)\in[-1,1]^2\times\{0,1\}:u+2z\le1,\ u^2\le s\},
 \qquad f=u-z/2,\qquad r=s.
\]
The equality $s=0$ forces $u=z=0$, so $v=0$. All native residuals are
nonnegative. Proposition~\ref{prop:mixture} therefore gives $D_\rho=0$
for every $\rho\ge0$.

This can also be checked directly. Put $t=\lambda+\rho$ and assume
$t\ge1/2$. The branch $z=1$ contains only $(-1,1,1)$ and has value
$t-3/2$. On $z=0$, minimization first sets $s=u^2$ and then
$u=-1/(2t)$, giving
\[
 L_\rho(\lambda)=\min\{t-3/2,-1/(4t)\}.
\]
This tends to zero as $\lambda\to+\infty$. For any finite $t$, however,
a sufficiently small $\eta>0$ gives the native point
$(-\eta,\eta^2,0)$ with value $-\eta+t\eta^2<0$.
Thus no finite multiplier attains the exact dual value. The source's reduction to multiplier zero invokes its Lemma~3,
which transfers attained exactness between fixed multipliers and does not
apply to this nonattained supremum. Thus the example demonstrates
nonattainment, not failure of dual-value exactness. This correction does
not affect the source's Theorem~14 under slice Slater conditions, nor the
attained lower-bound formulas in Section~\ref{sec:lower}.

For genuine failure of finite dual-value exactness without refined Slater,
the two-sided exceptional-atom construction in Section~\ref{sec:perturbation}
has an infinite optimized threshold at right-hand side zero.
